% - Návrh dotazů - jak jsou dotazy realizovány, jaká data a zdroje potřebují

\section{Návrh dotazů}
\label{sec:queries}

Tato sekce popisuje, jakým způsobem jsou analytické dotazy definované v~\customref{Subsekci}{subsec:requirements:queries} realizovány. Pro každý dotaz je uvedeno, z~jakých datových zdrojů čerpá a~které tabulky databáze využívá.

Všechny dotazy jsou striktně čtecí, parametrizované identifikátorem repozitáře a~pracují výhradně s~daty uloženými v~databázi \sectionref{sec:database}. Tím je zaručena nízká latence odpovědí i~nad velkým množstvím historických záznamů.

\subsection{Historická dostupnost dat}
\label{subsec:queries:historical_data}

Dotazy Q1, Q2 a~Q4 pracují s~historickým vývojem \PR a~issues. Aby bylo možné tyto dotazy spolehlivě zodpovědět, musí mít systém v~databázi kompletní sekvenci událostí pro každý \PR a~issue, včetně změn štítků a~stavů. Nestačí tedy ukládat pouze aktuální stav, je nutné uchovávat celou historii změn.

V~prostředí GitHub existují dva přístupy k~získání těchto dat:

\begin{description}
  \item[Webhooky] GitHub při každé relevantní události zašle \HTTP požadavek na předem registrovaný endpoint~\cite{github_webhooks}. Tento přístup je vhodný pro reaktivní systémy pracující v~reálném čase, ale má zásadní omezení: vyžaduje veřejně dostupný endpoint a~neposkytuje historická data. Nelze jím zpětně získat události, ke kterým došlo před registrací webhooku nebo v~době, kdy aplikace nebyla dostupná.
  \item[Timeline events \API] GitHub pro každý \PR nebo issue zpřístupňuje endpoint vracející úplný chronologický seznam všech zaznamenaných událostí~\cite{github_rest_api}. Tento přístup umožňuje jak počáteční naplnění historických dat, tak inkrementální doplňování nových událostí. Nevýhodou je vyšší počet \API požadavků při prvotní synchronizaci, avšak pro splnění požadavku na retroaktivní dostupnost dat je toto řešení nezbytné.
\end{description}

Systém proto používá Timeline events \API jako primární zdroj historických dat. Implementační detaily jsou popsány v~\sectionref{subsec:implementation:gotchas:timeline_vs_webhooks}.

\subsection{Realizace jednotlivých dotazů}
\label{subsec:queries:design}

\begin{description}
  \item[Q1 -- Historický stav \PR k~danému datu]
    Dotaz rekonstruuje stavy~\PR v~zadaném časovém rozmezí. Data jsou čerpána z~\REST \API GitHubu \subsectionref{subsec:github_api_rest}, konkrétně z~endpointu pro Timeline events, kde jsou zaznamenávány stavové štítky začínající prefixem \code{S-*} \sectionref{sec:rust_labels}. V~databázi dotaz využívá tabulky \code{issues}, \code{issue\_labels\_history} a~\code{issue\_event\_history}.

  \item[Q2 -- Počet \PR v~konkrétním stavu k~datu]
    Dotaz počítá~\PR v~daném stavu k~zadanému datu. Využívá tytéž datové zdroje jako Q1. Pracuje s~tabulkami \code{issues} a~\code{issue\_labels\_history}. Pro každý~\PR je potřeba určit jeho poslední známý stavový štítek před zadaným datem. Opakovaným voláním tohoto dotazu pro různá data lze na prezentační vrstvě sestavit časovou řadu vývoje počtu~\PR v~konkrétním stavu.

  \item[Q3 -- \PR čekající na zpracování]
    Na rozdíl od Q1 a~Q2 tento dotaz nepracuje s~historickými daty, ale s~aktuálním stavem. V~databázi vyhledává~\PR, jejichž poslední stavový štítek odpovídá jednomu ze tří čekacích stavů. Využívá tabulky \code{issues} a~\code{issue\_labels\_history}.

  \item[Q4 -- Události konkrétní issue k~datu]
    Data pro tento dotaz pocházejí z~Timeline \API GitHubu \subsectionref{subsec:github_api_rest}, kde jsou zaznamenány akce jako uzavření, znovuotevření nebo komentáře. V~databázi dotaz využívá tabulky \code{issues} a~\code{issue\_event\_history}.

  \item[Q5 -- Nejčastěji upravované soubory přispěvatelem]
    Data o~aktivitě nad soubory plynou z~lokálního git repozitáře \sectionref{sec:git}. Pro každý merge commit spojený s~\PR jsou extrahovány dotčené soubory a~uloženy do tabulky \code{file\_activity} spolu s~identifikátorem přispěvatele. Dotaz dále využívá tabulku \code{contributors} pro metadata přispěvatele.

  \item[Q6 -- Přispěvatelé, kteří upravili soubor nebo složku]
    Tento dotaz nabízí inverzní pohled k~Q5. Cesta je vyhledávána jako prefix, takže dotaz pokrývá jak konkrétní soubor, tak celý adresář. Data pocházejí ze všech tří hlavních zdrojů: lokálního git repozitáře \sectionref{sec:git}, \REST \API GitHubu \subsectionref{subsec:github_api_rest} a~Team \API \sectionref{sec:rust_team_api}. V~databázi využívá tabulky \code{file\_activity}, \code{contributors} a~\code{contributors\_teams}.

  \item[Q7 -- Soubory upravené členy týmu]
    Tento dotaz kombinuje data ze všech tří zdrojů. V~databázi vyžaduje propojení tabulek \code{file\_activity}, \code{contributors\_teams} a~\code{teams}. Výsledek lze vrátit jako plochý seznam, stromovou strukturu seskupenou do zvolené hloubky, nebo jako úplný adresářový strom. Tato flexibilita umožňuje prezentační vrstvě \subsectionref{subsec:architecture:presentation_layer} přizpůsobit vizualizaci jak detailním grafům na úrovni souborů, tak přehledovým heatmapám celých modulů.
\end{description}
